\documentclass[conference]{IEEEtran}
\IEEEoverridecommandlockouts

\usepackage{cite}
\usepackage{amsmath,amssymb,amsfonts}
\usepackage{algorithm}
\usepackage{algorithmic}
\usepackage{booktabs}
\usepackage{graphicx}
\usepackage{textcomp}
\usepackage{xcolor}
\usepackage{url}
\usepackage[hidelinks]{hyperref}
\usepackage{xspace}

\graphicspath{{figures/}}
\input{generated/metrics.tex}

\newcommand{\gpgomea}{GP-GOMEA\xspace}
\newcommand{\cpuoriginal}{\texttt{cpu\_original}}
\newcommand{\exactgpu}{\texttt{gpu\_exact\_gom}}
\newcommand{\batchgpu}{\texttt{gpu\_batch\_gom}}
\newcommand{\aq}{\operatorname{AQ}}
\newcommand{\mainbodyendpage}{?}
\AtEndDocument{%
  \typeout{====================================================}%
  \typeout{[PAGE BUDGET] Main body + references end on page \mainbodyendpage\space (limit: 6, Appendix excluded).}%
  \typeout{====================================================}%
}

\begin{document}
\bstctlcite{IEEEexample:BSTcontrol}

\title{Parallelizing GP-GOMEA Symbolic Regression with Exact and Batched GPU Fitness Evaluation\\
{\footnotesize Final Project --- Parallel Programming, Spring 2026}}

\author{%
\IEEEauthorblockN{Shao Chi Huang, Hao Chun Fu, and Po Hsuan Lai}
\IEEEauthorblockA{Parallel Programming Final Project}}

\maketitle

\begin{abstract}
Gene-pool Optimal Mixing Evolutionary Algorithms for genetic programming
(\gpgomea) repeatedly test small candidate modifications to symbolic-regression
trees. The fitness calculation is data parallel, but the accept--reject control
flow is sequential: launching one GPU evaluation for each candidate GOM move can
be slower than the CPU, while batching moves changes the search. This project
implements and evaluates two CUDA backends for \gpgomea. The exact GPU backend
preserves immediate GOM decisions and accelerates only the fitness evaluator;
the batched backend evaluates many candidate programs per launch and commits
updates after a local population batch. Expressions are serialized into postfix
token programs and interpreted by CUDA blocks over training samples. Across a
scaling study that varies dataset size, population size, and batch size, exact
GPU evaluation reaches a median \ExactMedianSpeedup$\times$ wall-clock speedup
over the CPU baseline, while the best batch configuration reaches
\BestBatchMedianSpeedup$\times$ and a maximum \BestBatchMaxSpeedup$\times$. A
100-generation study shows that batching does not cause a statistically
significant final fitness regression in this subset, but it changes trajectories
and must be treated as an approximate parallel algorithm.
\end{abstract}

\begin{IEEEkeywords}
parallel genetic programming, CUDA, symbolic regression, GP-GOMEA, GPU batching
\end{IEEEkeywords}

\vspace{0.5em}
\noindent\fbox{\parbox{0.97\columnwidth}{\textbf{Project Category:}\quad
$\boxtimes$~\textbf{A} (Self-chosen topic)\qquad
$\square$~\textbf{B} (AI-agent code optimization)}}

\section{Introduction}

Symbolic regression searches for an analytic expression that predicts a target
variable from input features. Genetic programming (GP) is a natural fit because
it evolves expression trees directly~\cite{koza1992genetic,polirefieldguide}, but
tree variation is expensive: every candidate expression must be evaluated over a
training set. \gpgomea improves the search by using linkage learning and
Gene-pool Optimal Mixing (GOM), which replaces linked tree positions with donor
material and keeps a change only if fitness does not worsen~\cite{virgolin2017scalable,virgolin2021improving}.

The parallel-programming challenge is that \gpgomea exposes two incompatible
forms of parallelism. Fitness evaluation over independent training samples is
highly data parallel, but GOM decisions are sequential: the next candidate
depends on whether previous moves were accepted. A naive GPU implementation
launches one small kernel per meaningful move and can lose to CPU launch
overhead. A batched implementation improves occupancy and amortizes launches,
but delayed commits alter the algorithm. This project asks where the exact GPU
crossover occurs, how much batching improves throughput, and how much search
distortion batching introduces.

The contribution is a CUDA implementation and evaluation of three backends:
the original CPU path, a control-flow exact GPU path, and a delayed-commit
batched GPU path. The design linearizes pointer-heavy expression trees into
postfix token programs, maps one candidate program to a CUDA block, parallelizes
across samples inside the block, and uses CUDA Graph replay, reusable pinned
buffers, capacity-managed device buffers, and host-side serialization caches to
reduce overhead.

\section{Background and Related Work}

\subsection{GP-GOMEA and the Baseline}

\gpgomea represents each individual with a fixed tree template. Some nodes may
be inactive, but every location has a stable index. A Family of Subsets (FOS)
describes linked locations that should be mixed together~\cite{thierens2013optimal}.
For each parent, GOM visits FOS subsets in random order, chooses a donor, copies
the selected operators, applies coefficient mutation, and evaluates the modified
tree if the expressed model changed. A lower MSE is better; a non-worsening move
is accepted immediately, and a worsening move is rolled back. This immediate
update visibility is part of the algorithm, not an implementation detail.

The reference symbolic-regression paper studies small-expression
\gpgomea~\cite{virgolin2021improving}. Its setup motivates our function set and
selected comparison to the fixed-population GP-GOMEA setting denoted \(G\) in
that paper, including the analytic quotient operator
\(\aq(a,b)=a/\sqrt{1+b^2}\)~\cite{ni2013analytic}. We do not rerun the full
reference-paper matrix; our main focus is parallel acceleration and the
algorithmic cost of batching.

\subsection{GPU Genetic Programming}

GPU GP systems commonly exploit data-parallel fitness evaluation using compiled
GPU programs, linearized programs, or interpreters~\cite{harding2007fast,langdon2008simd,langdon2010many}.
Later comparisons of kernel-per-individual, kernel-per-generation, and
interpreter designs show that no mapping is universally best because compilation,
launch, and execution overheads trade off differently~\cite{kim2017gpgpgpu}. Our
workload is finer-grained than conventional generational GP: candidate programs
are small, meaningful candidates appear throughout GOM, and exact accept--reject
decisions are sequential. The core question is therefore not simply ``can
fitness run on a GPU?'' but where exact GPU evaluation becomes worthwhile, and
how much search distortion is traded for throughput when candidate evaluations
are batched.

\section{Methodology}

\subsection{Program Representation and Kernel Mapping}

CUDA kernels never consume C++ \texttt{Node*} or virtual \texttt{Op*} objects.
Each active expression is serialized on the host into a compact postfix sequence
of \texttt{GpuToken} records. A token stores an opcode, an optional feature
index, and an optional constant. Batched evaluation stores programs in a
fixed-stride padded array and stores active lengths separately.

The GPU kernel assigns one candidate program to one CUDA block. Threads inside
the block evaluate the same postfix instruction sequence on different training
samples and reduce the per-sample losses into a candidate MSE. This mapping
avoids the dominant structural warp-divergence problem: lanes in a warp do not
interpret unrelated trees. Different tree lengths are handled across blocks,
where the GPU scheduler can overlap work from many programs.

\subsection{Algorithms}

Algorithm~\ref{alg:gom} shows the sequential GOM skeleton. The exact GPU backend
follows this control flow and replaces only the fitness call with immediate GPU
evaluation. Algorithm~\ref{alg:batch} shows the batched backend. It clones a
population batch into local candidates, evaluates the meaningful candidates from
one FOS step in a single GPU launch, and commits the local candidates only after
the batch finishes.

\begin{algorithm}[t]
\caption{Sequential GOM for one parent}
\label{alg:gom}
\footnotesize
\begin{algorithmic}[1]
\REQUIRE parent \(x\), population \(P\), FOS \(\mathcal{F}\)
\STATE \(x' \leftarrow clone(x)\), \(f \leftarrow fitness(x)\)
\FOR{each \(S \in shuffle(\mathcal{F})\)}
  \STATE \(d \leftarrow uniform(P)\)
  \STATE \((r,meaningful) \leftarrow MixAndMutate(x',d,S)\)
  \IF{\(meaningful\)}
    \STATE \(f' \leftarrow Fitness(x')\)
    \IF{\(f' \le f\)}
      \STATE discard rollback \(r\); \(f \leftarrow f'\)
    \ELSE
      \STATE \(Rollback(x',r,f)\)
    \ENDIF
  \ELSE
    \STATE discard \(r\)
  \ENDIF
\ENDFOR
\RETURN \(x'\)
\end{algorithmic}
\end{algorithm}

\begin{algorithm}[t]
\caption{Batched GPU GOM for one generation}
\label{alg:batch}
\footnotesize
\begin{algorithmic}[1]
\REQUIRE population \(P\), FOS \(\mathcal{F}\), batch size \(B\)
\FOR{each batch \(I\) of at most \(B\) population slots}
  \STATE \(C \leftarrow \{c_i = clone(P_i): i \in I\}\)
  \STATE \(O_i \leftarrow shuffle(\mathcal{F})\) for each \(c_i \in C\)
  \STATE \(f_i \leftarrow fitness(c_i)\) for each \(c_i \in C\)
  \FOR{\(s = 1,\ldots,|\mathcal{F}|\)}
    \STATE \(Q \leftarrow \emptyset\)
    \FOR{each candidate \(c_i \in C\)}
      \STATE \(S \leftarrow O_i[s]\), \(d \leftarrow uniform(P)\)
      \STATE \((r_i,m_i) \leftarrow MixAndMutate(c_i,d,S)\)
      \IF{\(m_i\)}
        \STATE \(Q \leftarrow Q \cup \{(i,r_i)\}\)
      \ELSE
        \STATE discard \(r_i\)
      \ENDIF
    \ENDFOR
    \STATE \(G \leftarrow GPUFitness(\{c_i \mid (i,r_i) \in Q\})\)
    \FOR{each \((i,r_i) \in Q\)}
      \IF{\(G_i \le f_i\)}
        \STATE accept move; \(f_i \leftarrow G_i\)
      \ELSE
        \STATE \(Rollback(c_i,r_i)\)
      \ENDIF
    \ENDFOR
  \ENDFOR
  \STATE commit each \(c_i \in C\) back to slot \(P_i\)
\ENDFOR
\end{algorithmic}
\end{algorithm}

\subsection{Optimization Techniques}

The optimizations target three bottlenecks. First, tree objects are converted to
compact postfix token programs before reaching CUDA. This removes host pointers
and virtual dispatch from device execution and lets each kernel interpret a
contiguous opcode stream. Second, the exact path is launch-bound, so it captures
the repeated copy--clear--kernel--copy sequence with CUDA Graph replay~\cite{nvidiacuda}, uses
pinned host buffers for transfers, keeps capacity-managed device buffers across
evaluations, and stores training data in single precision on the device. These
choices reduce per-candidate API, transfer, and allocation overhead while
preserving immediate GOM decisions. Third, the batch path uses a fixed-stride
token matrix with per-program lengths so many meaningful candidates can be
evaluated by one grid launch. Reused subtree vectors and serialization buffers
reduce host-side allocation during candidate collection. A persistent-kernel
prototype was tested but rejected because queue management and synchronization
overhead outweighed launch savings for this workload.

\section{Experimental Setup}

Experiments ran on an AMD Ryzen 9 7950X CPU, 61~GiB of memory, and two NVIDIA
RTX 4090 GPUs. CUDA code was built for compute capability 8.9. GPU timing uses
at most one process per physical GPU; CPU baselines are run sequentially to avoid
contention. All main suites use five search seeds, fixed seed-0 train/validation
/test splits, function set \(+,-,\times,\aq\), MSE fitness, and IMS disabled.

We measure wall-clock speedup \(S=t_{CPU}/t_{GPU}\) and evaluation-throughput
ratio separately, because batch search may evaluate a different number of
candidates. Correctness is checked at three levels: GPU batch versus GPU single
fitness, GPU versus CPU fitness with scale-aware floating-point tolerance, and
end-to-end validation/test reevaluation of archived elites. The exact GPU path
preserves GOM control flow but is not bitwise identical because device math and
parallel reductions can change ties. The batch GPU path is deliberately
approximate because population commits are delayed.

The acceleration suite contains \AccelerationRuns{} runs over four datasets
(Wine white, Tower, Airfoil, Dow chemical), sample scales 8, 16, and 32,
yielding 4,264--80,000 effective training samples per candidate evaluation,
population sizes 512, 1024, 2048, 4096, and 8192, height 3, one generation,
exact GPU, and batch sizes 64, 128, 256, 512, and 1024. Sample scaling repeats
training samples to stress throughput and is not used for statistical accuracy
claims. A separate 100-generation fitness suite contains \FitnessRuns{} runs and
\GenerationRecords{} generation records over four unscaled datasets, tree
heights 3, 4, and 5, population 1024, and the same batch sizes. We also run a
selected 20-run subset of the reference paper's fixed-population GP-GOMEA
setting, denoted \(G\) in that paper; the full 30-repetition, 1000-second
reference-paper matrix is not rerun.

To test the kernel-mapping assumption directly, we also run a synthetic CUDA
ablation outside the GOM loop. It compares the production-style
one-program-per-block mapping against a lane-mixed kernel where lanes in a warp
interpret different programs. Both evaluate 1,024 postfix programs over 32,768
samples; the heterogeneous case varies program length among 7, 15, 23, and 31
tokens and reports active token-samples per second.

\section{Results}

\subsection{Acceleration and Scaling}

Table~\ref{tab:accel} summarizes the main acceleration suite. Exact GPU
evaluation is faster than CPU in \ExactFasterFraction{} of paired runs and has
median speedup \ExactMedianSpeedup$\times$ with 95\% bootstrap confidence
interval \ExactSpeedupCI$\times$. Batching is consistently faster in this matrix:
\BestBatchVariant{} reaches median \BestBatchMedianSpeedup$\times$ and maximum
\BestBatchMaxSpeedup$\times$.

\begin{table}[t]
\centering
\caption{Acceleration-suite speedups relative to \cpuoriginal. Each row
summarizes 300 paired one-generation runs with 4,264--80,000 effective training
samples per candidate evaluation.}
\label{tab:accel}
\resizebox{\columnwidth}{!}{\input{generated/table_acceleration_summary.tex}}
\end{table}

\begin{figure}[t]
\centering
\includegraphics[width=\columnwidth]{speedup_vs_population_x32.png}
\caption{Median speedup at sample scale 32, corresponding to 17,056--80,000
effective training samples per candidate evaluation. Larger populations expose
more GOM candidates and amortize GPU setup.}
\label{fig:population}
\end{figure}

Figure~\ref{fig:population} shows the main crossover behavior. Exact GPU
benefits from larger data and population sizes but remains sensitive to launch
and synchronization overhead. Batch GPU reaches much higher throughput because a
single launch carries many candidate programs.

Table~\ref{tab:batchscale} isolates requested batch size on one large workload.
The requested size is an upper bound: actual GOM batches are smaller when fewer
meaningful candidates are available. Once the actual batch reaches tens of
programs per launch, wall-clock speedup jumps from single digits to roughly
46--51$\times$ on this case.

\begin{table}[t]
\centering
\caption{Batch-size scaling on Wine white, sample scale 32, population 4096.
This setting evaluates 78,368 effective training samples per candidate.}
\label{tab:batchscale}
\resizebox{\columnwidth}{!}{\input{generated/table_batch_scaling.tex}}
\end{table}

\subsection{Profiling and Bottleneck Explanation}

Nsight Systems profiling on Wine white sample scale 32, population 4096, height
3, seed 0 evaluates 78,368 effective training samples per candidate and exposes
the root cause. Table~\ref{tab:profile} shows that exact GPU uses
about one candidate per effective CUDA Graph launch, causing CUDA API time to
dominate. Batch GPU evaluates hundreds of programs per launch, reducing host API
time per candidate even though total kernel time is higher.

\begin{table}[t]
\centering
\caption{Representative Nsight Systems profile.}
\label{tab:profile}
\resizebox{\columnwidth}{!}{\input{generated/table_profiling.tex}}
\end{table}

The synthetic mapping ablation gives a separate check on warp-level control
flow. With variable-length programs, median active token-sample throughput over
ten repeats is 0.97$\times$ the uniform-length case under one-program-per-block
mapping, but only 0.60$\times$ under lane-mixed mapping. This supports the
choice to keep one expression tree inside a block rather than mixing unrelated
trees within a warp.

A before--after ablation on the representative workload shows that these
engineering changes reduce exact-GPU wall time by \ExactOptimizationSpeedup$\times$
and batch-GPU wall time by \BatchOptimizationSpeedup$\times$, confirming that
launch amortization and buffer reuse are the main optimization targets.

\subsection{Fitness Distortion}

Batching changes the search trajectory because donor visibility is delayed, so
speed alone is insufficient. Table~\ref{tab:distortion} compares final test NMSE
after 100 generations against exact GPU. No batch size shows a significant
paired final test-NMSE difference after Holm correction at the 0.05 level in this
subset. However, the median ratios are not identical, and the win/loss counts
confirm that batching is an algorithmic variant rather than a transparent
implementation optimization.
Figure~\ref{fig:fitness} visualizes the same point over generations: quality is
often close by the end, but the methods follow different search paths.

\begin{table}[t]
\centering
\caption{Final batch distortion relative to exact GPU. Lower test ratio is better.}
\label{tab:distortion}
\resizebox{\columnwidth}{!}{\input{generated/table_fitness_distortion.tex}}
\end{table}

\begin{figure}[!t]
\centering
\includegraphics[width=\columnwidth]{fitness_evolution_by_generation.png}
\caption{Median best test NMSE over 100 generations for height 4. Batched
variants can reach useful quality quickly, but do not trace the same trajectory
as exact GOM.}
\label{fig:fitness}
\end{figure}

\begin{figure}[!t]
\centering
\includegraphics[width=\columnwidth]{quality_speed_tradeoff.png}
\caption{Batch size trades throughput against generation-matched quality.}
\label{fig:quality}
\end{figure}

\section{Discussion and Analysis}

The results support three conclusions relevant to parallel programming. First,
fitness evaluation is not automatically GPU-friendly. The arithmetic work is
parallel, but exact GOM exposes it in very small sequential pieces. CUDA Graphs
and pinned buffers reduce overhead, yet exact speedup depends on enough training
samples and enough candidates to amortize each decision. This explains why exact
GPU reaches higher speedup at sample scale 32 but is not always dominant on
small cases.

Second, batching is the main source of throughput because it increases both
sample-level and program-level parallelism. The selected block mapping addresses
tree-induced warp divergence by keeping one tree per block: lanes in a warp
execute the same postfix operations over different samples. Different tree
shapes become inter-block imbalance rather than intra-warp branch divergence,
which matches the mapping ablation reported above. The profiling result of 334
programs per launch for batch GPU versus roughly one program per launch for
exact GPU is the most direct evidence for launch amortization.

Third, batching has a scientific cost. Delayed commits change which donors are
visible and therefore change the sequence of accepted GOM moves. This is why the
paper reports both speedup and final fitness distortion. For this experiment,
batching provides large speedups without a statistically significant final
quality loss, but this should not be generalized without the full paper matrix,
more repetitions, and additional datasets.

The main limitations are scope and hardware dependence. Sample scaling increases
per-candidate work for throughput stress tests but does not add new statistical
information. The selected \(G\)-setting rerun is only an external reference
check, not a protocol-equivalent reproduction of the reference paper's full
30-repetition, 1000-second matrix. Finally, the crossover point between CPU,
exact GPU, and batched GPU depends on GPU launch latency, memory bandwidth, CPU
scalar performance, and available population-level parallelism.

\section{Conclusion}

This project parallelizes \gpgomea symbolic-regression fitness evaluation with
two CUDA execution modes. Exact GPU evaluation preserves GOM update semantics and
achieves a median \ExactMedianSpeedup$\times$ speedup on the scaling study.
Batched GPU evaluation exposes much more parallel work and reaches median
\BestBatchMedianSpeedup$\times$ speedup, but it is an approximate delayed-commit
algorithm. The key lesson is that parallelizing an evolutionary algorithm
requires reporting both hardware efficiency and algorithmic distortion: the
fastest implementation is not necessarily the same algorithm.

\bibliographystyle{IEEEtran}
\bibliography{references}

\xdef\mainbodyendpage{\thepage}

\end{document}
